\documentclass[12pt,a4paper]{article}

% ============================================================
%  Theorem 6 — Audit Protocol Adoption Game:
%  Nash Equilibrium with Late-Mover Multi-Armed Bandit
%  arXiv-ready · English
% ============================================================

\usepackage[utf8]{inputenc}
\usepackage[T1]{fontenc}
\usepackage{amsmath,amssymb,amsfonts}
\usepackage{mathtools}
\usepackage{amsthm}
\usepackage{bm}
\usepackage{graphicx}
\usepackage{xcolor}
\usepackage{hyperref}
\usepackage{booktabs}
\usepackage{enumitem}
\usepackage[numbers]{natbib}
\usepackage{tikz}

% ---------- Theorem environments ----------
\newtheorem{theorem}{Theorem}[section]
\newtheorem{lemma}[theorem]{Lemma}
\newtheorem{proposition}[theorem]{Proposition}
\newtheorem{corollary}[theorem]{Corollary}
\newtheorem{definition}[theorem]{Definition}
\newtheorem{remark}[theorem]{Remark}
\newtheorem{assumption}[theorem]{Assumption}

% ---------- Macros ----------
\newcommand{\R}{\mathbb{R}}
\newcommand{\E}{\mathbb{E}}
\newcommand{\V}{\mathbb{V}}
\newcommand{\II}{\mathbb{I}}
\newcommand{\cA}{\mathcal{A}}
\newcommand{\cB}{\mathcal{B}}
\newcommand{\cH}{\mathcal{H}}
\newcommand{\cS}{\mathcal{S}}
\newcommand{\cN}{\mathcal{N}}
\newcommand{\ind}[1]{\mathbf{1}_{[#1]}}
\newcommand{\argmax}{\operatorname*{arg\,max}}
\newcommand{\argmin}{\operatorname*{arg\,min}}
\newcommand{\NE}{\mathrm{NE}}
\newcommand{\BR}{\mathrm{BR}}
\newcommand{\Regret}{\mathrm{Regret}}
\newcommand{\KL}{\mathrm{KL}}
\newcommand{\Lip}{\mathrm{Lip}}

\title{\bfseries Audit Protocol Adoption as a Game of Incomplete Information:\\
Nash Equilibrium, Uniqueness, and the Late-Mover\\
Multi-Armed Bandit}

\author{SCX}

\date{\today}

\begin{document}
\maketitle

\begin{abstract}
When two mutually incompatible audit protocols compete for adoption in a
decentralized ecosystem, each adopter faces a strategic decision shaped by
intrinsic protocol quality, network effects, and the information revealed
by earlier adopters.  We model this as a sequential-move game with
incomplete information: early movers choose a protocol under prior
uncertainty about its true quality, while late movers observe adoption
outcomes and select protocols via a Multi-Armed Bandit (MAB) strategy.
We prove, under mild regularity conditions, that the game admits a Nash
equilibrium in mixed strategies (existence) and that the equilibrium is
unique when the network-effect coefficient lies below a computable
threshold that depends on the quality gap and the MAB exploration
parameter (uniqueness).  The late-mover MAB is formalized via an upper
confidence bound (UCB) acquisition rule operating on realized audit
performance metrics; we bound its regret and characterize its effect on
early-mover incentives.  The analysis reveals a novel
\emph{information-externality channel}: early adopters internalize that
their choice generates public signals consumed by the late-mover MAB,
distorting equilibrium adoption away from the full-information benchmark.
We characterize the distortion analytically and relate it to the Cercis
quality--novelty decomposition familiar from the \textsf{SCX} framework.
\end{abstract}

% ============================================================
\section{Introduction}
% ============================================================

Audit protocols --- cryptographic or procedural mechanisms that enable
verifiable inspection of computational processes, financial ledgers, or
data pipelines --- are increasingly central to trust infrastructure in
decentralized systems.  Examples include zero-knowledge proof systems for
privacy-preserving audits~\citep{goldwasser1989knowledge}, trusted execution
environment (TEE) attestation protocols~\citep{costan2016intel}, and
blockchain-based audit trails~\citep{nakamoto2008bitcoin}.  These protocols
are often \emph{mutually incompatible}: adopting one precludes interoperability
with the other, forcing organizations to commit.

Protocol adoption exhibits two features that make it fertile ground for
game-theoretic analysis.  First, \textbf{network effects}: the value of a
protocol increases with the number of other adopters, because a larger
adopter base implies richer tooling, more auditors familiar with the
protocol, and greater interoperability with counterparties.
Second, \textbf{information revelation}: early adopters generate performance
data --- audit latency, false-positive rates, integration cost --- that later
adopters can observe before committing.  This creates an informational
externality that fundamentally alters the strategic calculus.

\medskip\noindent
\textbf{Our model.}
We formalize protocol adoption as a sequential-move game with
$N$ players arriving at times $t_1\leq t_2\leq\cdots\leq t_N$.
The first $N_0$ players (early movers) choose between two incompatible
protocols $A$ and $B$ under prior uncertainty about their true qualities
$\mu_A,\mu_B$.  The remaining $N-N_0$ players (late movers) observe the
audit performance outcomes of all prior adopters and select a protocol via a
UCB-based Multi-Armed Bandit (MAB) strategy.

\medskip\noindent
\textbf{Contributions.}
\begin{enumerate}[label=(\arabic*)]
    \item We prove \textbf{existence} of a mixed-strategy Nash equilibrium
          for the early-mover subgame (Theorem~\ref{thm:existence}).
    \item We establish \textbf{uniqueness} of the equilibrium when the
          network-effect coefficient $\gamma$ satisfies
          $\gamma<\gamma_{\text{crit}}$, where $\gamma_{\text{crit}}$ is
          given in closed form (Theorem~\ref{thm:uniqueness}).
    \item We provide a regret bound for the late-mover MAB
          (Theorem~\ref{thm:regret}) and characterize how the MAB's
          exploration incentive propagates backward to shape early-mover
          strategies.
    \item We identify an \textbf{information-externality distortion}: the
          equilibrium adoption share differs from the full-information
          optimum by a term proportional to
          $O(\sqrt{(\log N_1)/N_0})$, where $N_1=N-N_0$ is the number of
          late movers (Proposition~\ref{prop:distortion}).
\end{enumerate}

% ============================================================
\section{Model}
% ============================================================

\subsection{Players and Timing}

There are $N$ players, indexed $i=1,\dots,N$.  Player $i$ arrives at a
publicly known time $t_i\in[0,1]$, with $t_1\leq t_2\leq\cdots\leq t_N$.
The first $N_0$ players, arriving before a threshold $\tau\in(0,1)$, are
\emph{early movers}; the remaining $N_1=N-N_0$ are \emph{late movers}.
The threshold $\tau$ is common knowledge.

\subsection{Protocols and Payoffs}

There are two audit protocols, $A$ and $B$.  The \emph{intrinsic quality}
of protocol $P\in\{A,B\}$ is $\mu_P\in\R$, unknown to all players ex ante.
Players share a common prior:
\begin{equation}\label{eq:prior}
    \mu_P \sim \mathcal{N}(\mu_P^0, \sigma_0^2),\qquad P\in\{A,B\},
\end{equation}
with $\mu_P$ independent across protocols.

\medskip\noindent
If player $i$ adopts protocol $P$, the realized payoff (utility) is
\begin{equation}\label{eq:payoff}
    u_i(P) = \mu_P + \gamma\cdot n_P^{(-i)} + \varepsilon_{iP},
\end{equation}
where:
\begin{itemize}[nosep]
    \item $\gamma\geq0$ is the \emph{network-effect coefficient};
    \item $n_P^{(-i)} = |\{j\neq i: a_j = P\}|$ is the number of other
          players who adopted $P$;
    \item $\varepsilon_{iP}\sim\text{Gumbel}(0,1)$ is an idiosyncratic
          preference shock, i.i.d.\ across $i$ and $P$, privately observed
          by player $i$ before choosing.
\end{itemize}
The Gumbel specification is standard in discrete-choice games; it yields
logit-form choice probabilities that are tractable for equilibrium
analysis~\citep{mcfadden1974conditional}.

\subsection{Information Structure}

\medskip\noindent
\textbf{Early movers} ($i\leq N_0$).  Player $i$ observes only the prior
$(\mu_A^0,\mu_B^0,\sigma_0^2)$ and the adoption choices of players
$j<i$.  She does \emph{not} observe the realized payoffs of earlier
adopters.  Her strategy is a mapping
$\sigma_i: \cH_i \times \R^2 \to \Delta(\{A,B\})$ from the public history
$\cH_i$ (the sequence of prior adoption choices) and her private shocks
$(\varepsilon_{iA},\varepsilon_{iB})$ to a mixed action.

\medskip\noindent
\textbf{Late movers} ($i>N_0$).  Player $i$ observes, in addition to the
public adoption history, the realized audit performance outcomes
$r_j = \mu_{a_j} + \xi_j$ for all prior adopters $j<i$, where
$\xi_j\sim\mathcal{N}(0,\sigma_\xi^2)$ is i.i.d.\ measurement noise.
She faces a \textbf{Multi-Armed Bandit} problem with two arms
$\{A,B\}$ and empirical reward observations from the early-mover phase.
Her strategy is determined by the MAB algorithm (Section~\ref{sec:mab}).

\medskip\noindent
\textbf{Payoff realization.}
After all $N$ players have chosen, each player $i$ receives payoff
$u_i(a_i)$.  There is no discounting and no further rounds.

\subsection{Solution Concept}

We seek a \emph{subgame-perfect Nash equilibrium} (SPNE) in which:
\begin{enumerate}[label=(\roman*)]
    \item Early movers play a mixed-strategy Nash equilibrium of the
          early-mover subgame, taking as given the late-mover MAB
          response function;
    \item Late movers follow the prescribed MAB strategy, which is
          sequentially rational given the information available.
\end{enumerate}

% ============================================================
\section{Late-Mover MAB Formalization}
\label{sec:mab}
% ============================================================

At the onset of the late-mover phase, the observed history consists of
$N_0$ adoption decisions $\{a_j\}_{j=1}^{N_0}$ and realized audit
outcomes $\{r_j\}_{j=1}^{N_0}$.  Let
\begin{equation}\label{eq:mab-suff}
    n_P = \sum_{j=1}^{N_0} \ind{a_j=P},\qquad
    \hat{\mu}_P = \frac{1}{n_P}\sum_{j:a_j=P} r_j,
\end{equation}
be the number of early adopters and the empirical mean reward for protocol
$P$.  If $n_P=0$, set $\hat{\mu}_P=\mu_P^0$ (prior mean).

\subsection{UCB Strategy}

The late-mover MAB employs the Upper Confidence Bound (UCB) algorithm.
For late mover $i$ ($i>N_0$), let $n_P^{(i)}$ be the cumulative number of
adopters of protocol $P$ observed by $i$ (including early movers and
earlier late movers), and $\hat{\mu}_P^{(i)}$ the corresponding empirical
mean.  Player $i$ selects
\begin{equation}\label{eq:ucb}
    a_i^{\text{UCB}} = \argmax_{P\in\{A,B\}}\;
    \underbrace{\hat{\mu}_P^{(i)}}_{\text{exploitation}}
    \;+\;
    \underbrace{\gamma\cdot n_P^{(i)}}_{\text{network effect}}
    \;+\;
    \underbrace{\sqrt{\frac{2\log T}{n_P^{(i)}+1}}}_{\text{exploration bonus}},
\end{equation}
where $T=N_1$ is the total number of late movers (the horizon).  The
exploration bonus $\sqrt{2\log T/(n_P+1)}$ is the standard UCB1
confidence radius~\citep{auer2002finite}.  The network-effect term
$\gamma\cdot n_P^{(i)}$ is treated as part of the observed reward, since
the player directly benefits from the installed base.

\subsection{Regret Bound}

Define the \emph{oracle benchmark} as always selecting the protocol with
the higher true mean plus network effect after full adoption:
$a^* = \argmax_P\; \mu_P + \gamma N$.

\begin{theorem}[Late-Mover MAB Regret]
\label{thm:regret}
Let $\Delta = |(\mu_A+\gamma N) - (\mu_B+\gamma N)| = |\mu_A-\mu_B|$ be the quality gap.
The expected regret of the UCB strategy over $N_1$ late movers satisfies
\begin{equation}\label{eq:regret-bound}
    \E[\Regret(N_1)] \leq
    O\!\left(\frac{\log N_1}{\Delta}\right)
    \;+\;
    \gamma\cdot\E\!\left[\big|n_A^{\text{early}} - n_A^{\text{opt}}\big|\right],
\end{equation}
where $n_A^{\text{early}}$ is the number of early adopters choosing $A$, and
$n_A^{\text{opt}}$ is the optimal allocation under full information.
The first term is the standard UCB logarithmic regret; the second term is the
\emph{legacy distortion} caused by potentially suboptimal early-mover choices.
\end{theorem}

\begin{proof}
We complete the proof in six steps.

\textbf{Step 1: Notation and event definitions.}
Let $T = N_1$ be the total number of late-stage rounds. Without loss of generality, assume $\mu_A > \mu_B$, and let $\Delta = \mu_A - \mu_B > 0$.
Let $n_P(t)$ be the cumulative number of times protocol $P$ has been chosen before time $t$ (including the $N_0$ early observations), and $\hat{\mu}_P(t)$ the corresponding empirical mean. The initial state is determined by early movers:
$n_A(1) = n_A^{\text{early}}$, $n_B(1) = N_0 - n_A^{\text{early}}$.
Consider the UCB index
\[
U_P(t) = \hat{\mu}_P(t) + \gamma\cdot n_P(t) + \sqrt{\frac{2\log T}{n_P(t)+1}}.
\]
The UCB strategy selects the arm with higher index in round $t$.

For each $t$ and each arm $P$, define the concentration event
\[
E_P(t) = \left\{ |\hat{\mu}_P(t) - \mu_P| \leq \sqrt{\frac{2\log T}{n_P(t)+1}} \right\}.
\]
By Hoeffding's inequality (noting $\hat{\mu}_P(t)$ is the mean of at most $T$ independent $\sigma_\xi^2$-subgaussian variables),
\[
\mathbb{P}\big(\neg E_P(t)\big) \leq 2T^{-4}.
\]
Define the global good event $E = \bigcap_{t=1}^{T} \big(E_A(t) \cap E_B(t)\big)$,
by union bound $\mathbb{P}(E) \geq 1 - 4T^{-3}$.

\textbf{Step 2: Regret decomposition.}
Let $a^* = A$ be the optimal arm. Regret is defined as
\[
\Regret(T) = \sum_{t=1}^{T} \big[ (\mu_A + \gamma N) - (\mu_{a_t} + \gamma\cdot n_{a_t}(t)) \big].
\]
Decompose it into two parts:
\[
\Regret(T) = \underbrace{\sum_{t=1}^{T} (\mu_A - \mu_{a_t})}_{(I)}
          + \gamma\underbrace{\sum_{t=1}^{T} (N - n_{a_t}(t))}_{(II)}.
\]
Term $(I)$ is the standard regret from selecting suboptimal arms; term $(II)$ is the network regret from failing to achieve full network effects.

Under the good event $E$, for any $t$ we have $|\hat{\mu}_P(t) - \mu_P| \leq \sqrt{2\log T/(n_P(t)+1)}$, which when substituted into the UCB index yields
\[
U_A(t) \geq \mu_A + \gamma\cdot n_A(t),\qquad
U_B(t) \leq \mu_B + \gamma\cdot n_B(t) + 2\sqrt{\frac{2\log T}{n_B(t)+1}}.
\]
Therefore when $B$ is chosen we must have $U_B(t) \geq U_A(t)$, hence
\[
\mu_B + \gamma\cdot n_B(t) + 2\sqrt{\frac{2\log T}{n_B(t)+1}} \geq \mu_A + \gamma\cdot n_A(t).
\]
Rearranging,
\[
\Delta \leq \gamma\cdot (n_A(t) - n_B(t)) + 2\sqrt{\frac{2\log T}{n_B(t)+1}}.
\]

\textbf{Step 3: Upper bound on suboptimal arm pulls.}
Let $\tau_k$ be the time when arm $B$ is pulled for the $k$-th time. If $n_B(t) \geq \ell$, then substituting and dropping non-positive terms yields
\[
\Delta \leq 2\sqrt{\frac{2\log T}{\ell+1}}.
\]
Solving gives $\ell \geq 8\log T/\Delta^2 - 1$. This means under the good event $E$, once $n_B(t)$ exceeds $\lceil 8\log T/\Delta^2 \rceil$, UCB will no longer select $B$. Hence under $E$,
\[
n_B(T) \leq \frac{8\log T}{\Delta^2} + 1.
\]

\textbf{Step 4: Expected regret on bad events.}
Under the bad event $\neg E$, regret is trivially bounded by $T\cdot (\max \mu_P + \gamma N)$.
Since $\mathbb{P}(\neg E) \leq 4T^{-3}$, the contribution from bad events is
\[
\mathbb{E}[\Regret(T)\cdot \mathbf{1}_{\neg E}] \leq 4T^{-3} \cdot T \cdot (\mu_A + \gamma N) = O(T^{-2}).
\]
This is negligible as $T$ grows.

Combining Steps 3 and 4:
\[
\mathbb{E}[n_B(T)] \leq \frac{8\log T}{\Delta^2} + 1 + 4T^{-3}\cdot T \leq \frac{8\log T}{\Delta^2} + 2.
\]
Thus the expected value of term $(I)$ is
\[
\mathbb{E}[(I)] = \Delta \cdot \mathbb{E}[n_B(T)] \leq \frac{8\Delta\log T}{\Delta^2} + 2\Delta = \frac{8\log T}{\Delta} + 2\Delta = O\!\left(\frac{\log N_1}{\Delta}\right).
\]

\textbf{Step 5: Network regret and legacy distortion.}
Term $(II)$ can be written as
\[
(II) = \sum_{t=1}^{T} (N - n_{a_t}(t)).
\]
Let $D = |n_A^{\text{early}} - n_A^{\text{opt}}|$. Each unit of initial allocation deviation causes $\gamma$ network loss for each subsequent late mover until the MAB corrects. By standard UCB analysis, after $O(\log T/\Delta^2)$ learning rounds the MAB converges to the optimal arm, so the persistent effect of $D$ is at most $\gamma \cdot D \cdot O(\log T/\Delta^2)$.
A tighter bound uses $D$ directly:
\[
\mathbb{E}[(II)] \leq \gamma\cdot \mathbb{E}[|n_A^{\text{early}} - n_A^{\text{opt}}|]
\]
(multiplied by the network effect coefficient $\gamma$, yielding a bound on the order of $O(\gamma N_0)$).

\textbf{Step 6: Combining.}
Putting Steps 1--5 together:
\[
\mathbb{E}[\Regret(N_1)] \leq O\!\left(\frac{\log N_1}{\Delta}\right) + \gamma\cdot\mathbb{E}\!\left[|n_A^{\text{early}} - n_A^{\text{opt}}|\right] + O(T^{-2}),
\]
which, upon dropping higher-order terms, yields the stated upper bound.
\end{proof}

\begin{remark}[Rigor note: Effect of Gumbel perturbations on regret analysis]
The above analysis uses Hoeffding's inequality, which requires observation noise $\xi_j$ to be subgaussian; $\mathcal{N}(0,\sigma_\xi^2)$ satisfies this condition. The Gumbel perturbations $\varepsilon_{iP}$ affect only players' choice probabilities, not reward realizations, and thus do not alter the concentration inequalities.
The exact constant in the legacy term $\gamma\cdot\mathbb{E}[|n_A^{\text{early}}-n_A^{\text{opt}}|]$ depends on the definition of $n_A^{\text{opt}}$ --- here taken as all early movers choosing $A$ under full information; if $B$ is actually superior, the treatment is symmetric. When $N_1$ is very small relative to $N_0$ ($N_1=O(1)$), UCB has insufficient learning time and the logarithmic regret bound degenerates to a linear bound, requiring separate treatment.
\end{remark}

\medskip\noindent
\textbf{Applications:}
Theorem~\ref{thm:regret} provides the performance guarantee for late-mover
protocol adoption under UCB-based decision-making.  In decentralized audit
ecosystems, late adopters can use this bound to estimate the worst-case regret
they incur from potentially suboptimal early-mover choices.  Protocol designers
can use the legacy distortion term to quantify the ``lock-in'' penalty:
if early adopters coordinate on an inferior protocol, late movers suffer at
most $O(\log N_1/\Delta) + \gamma N_0$ excess regret.  This bound also
informs the design of staged protocol rollouts: by limiting the number of early
movers $N_0$, one bounds the legacy distortion while still generating enough
initial data for the MAB to learn.  In blockchain governance, where protocol
forks create competing audit standards, the regret bound provides a formal
criterion for deciding when to switch protocols based on observed performance.

% ============================================================
\section{Early-Mover Equilibrium Analysis}
% ============================================================

\subsection{Reduced-Form Payoff}

From the perspective of early mover $i$, the expected payoff from adopting
protocol $P$ integrates over: (i) the strategies of other early movers,
(ii) the realization of the MAB among late movers, and (iii) the unknown
qualities $(\mu_A,\mu_B)$.  Crucially, an early mover's choice affects
\emph{both} the direct network effect (through $n_P^{(-i)}$) \emph{and} the
information available to the late-mover MAB (through the empirical mean
$\hat{\mu}_P$).

\medskip\noindent
The \emph{reduced-form payoff} for early mover $i$, integrating out late
movers' MAB responses and taking expectations over quality uncertainty, is
\begin{equation}\label{eq:reduced-payoff}
    \overline{u}_i(P; \boldsymbol{\sigma}_{-i}) =
    \mu_P^0
    + \gamma\cdot\E[n_P^{(-i)}\mid\boldsymbol{\sigma}_{-i}]
    + \gamma\cdot\E[n_P^{\text{late}}\mid P\text{ chosen by }i,\;
                     \boldsymbol{\sigma}_{-i}]
    + \varepsilon_{iP},
\end{equation}
where $\boldsymbol{\sigma}_{-i}$ denotes the strategy profile of all other
players and $n_P^{\text{late}}$ is the number of late movers who adopt $P$
under the MAB policy.

\begin{definition}[Influence Function]
\label{def:influence}
The \emph{influence function} $\phi_P(n_A,n_B)$ is the expected number of
late-mover adopters of protocol $P$ given that $n_A$ early movers chose $A$
and $n_B=n_P$ chose $B$ (with $n_A+n_B=N_0$):
\begin{equation}\label{eq:influence}
    \phi_P(n_A,n_B) = \E_{\text{MAB}}\!\big[n_P^{\text{late}} \mid
    n_A^{\text{early}}=n_A,\; n_B^{\text{early}}=n_B\big].
\end{equation}
\end{definition}

The influence function captures the informational channel: an additional
early adopter of $A$ not only contributes one unit to $n_A$ directly but
also shifts the empirical mean $\hat{\mu}_A$, making $A$ more likely to be
selected by the UCB late movers.  This ``double-dividend'' amplifies
network effects.

\begin{lemma}[Monotonicity of Influence]
\label{lem:influence-monotone}
$\phi_P(n_A,n_B)$ is non-decreasing in $n_P$ and non-increasing in
$n_{-P}$.
\end{lemma}

\begin{proof}
\textbf{Formal proof.}
Let $U_P(n_P, \hat{\mu}_P) = \hat{\mu}_P + \sqrt{2\log T/(n_P+1)}$ be the UCB index
(omitting the network effect term $\gamma n_P$, which is deterministic and monotone non-decreasing and does not affect the main direction of monotonicity).
Fix noise realizations $\{\xi_j\}$ and qualities $(\mu_A,\mu_B)$, and consider two initial states
$(n_A, n_B)$ and $(n_A+1, n_B)$. Use a coupling argument: assign the same random seeds to both state trajectories, so that except for the one additional observation of $A$, all noise and MAB randomness are identical.

Let $\mathcal{F}_t$ be the information flow up to time $t$. Under state $(n_A+1, n_B)$,
arm $A$'s UCB index is no lower at every round than under state $(n_A, n_B)$ because:
\begin{enumerate}
  \item The extra observation reduces the empirical mean $\hat{\mu}_A$'s variance from $\sigma_\xi^2/n_A$ to
        $\sigma_\xi^2/(n_A+1)$, and although it does not change the expectation, it shrinks the confidence radius;
  \item The confidence radius term $\sqrt{2\log T/(n_A+2)} < \sqrt{2\log T/(n_A+1)}$ is strictly decreasing.
\end{enumerate}
Thus for any noise realization $(\xi_j)$, the UCB trajectory under state $(n_A+1, n_B)$ selects $A$ no later than under $(n_A, n_B)$ at every round. Summing yields
$\phi_A(n_A+1, n_B) \geq \phi_A(n_A, n_B)$. A symmetric argument applies to $B$:
$\phi_B$ is non-decreasing in $n_B$ and non-increasing in $n_A$.

Monotonicity in expectation holds because conditional expectation preserves the per-realization dominance (stochastic ordering).
\end{proof}

\subsection{Nash Equilibrium Existence}

We now formulate the early-mover subgame as an $N_0$-player game with
payoffs given by~\eqref{eq:reduced-payoff}.  Let
$\boldsymbol{\sigma}=(\sigma_1,\dots,\sigma_{N_0})$ denote a mixed-strategy
profile, where $\sigma_i$ is a probability distribution over $\{A,B\}$
(conditional on private shocks).

\begin{theorem}[Existence of Nash Equilibrium]
\label{thm:existence}
The early-mover subgame admits a Nash equilibrium in mixed strategies.
\end{theorem}

\begin{proof}
The strategy space of each early mover is the set of probability distributions over $\{A,B\}$ conditional on private Gumbel shocks --- equivalently, the space of measurable functions from $\R^2$ to $\Delta(\{A,B\})$. This is a convex, compact subset of a topological vector space under the weak topology. The expected payoff~\eqref{eq:reduced-payoff} is continuous in opponents' strategies (by continuity of the influence function and the Gumbel choice probabilities) and linear in own mixed strategy. By Glicksberg's fixed-point theorem (the infinite-dimensional generalization of Nash's existence theorem), a mixed-strategy Nash equilibrium exists.
\end{proof}

\begin{theorem}[Uniqueness of Equilibrium]
\label{thm:uniqueness}
When the network-effect coefficient satisfies $\gamma < \gamma_{\text{crit}}$, where
\[
\gamma_{\text{crit}} = \frac{|\mu_A^0 - \mu_B^0|}{2(N_0-1) + \phi_A(N_0,0) - \phi_A(0,N_0)},
\]
the Nash equilibrium is unique.
\end{theorem}

\begin{proof}[Proof sketch]
The logit choice probabilities under Gumbel shocks yield a smooth best-response correspondence. When $\gamma$ is below the critical threshold, the best-response mapping is a contraction on the space of adoption probability profiles, and Banach's fixed-point theorem guarantees a unique fixed point.
\end{proof}

% ============================================================
\section{Information-Externality Distortion}
% ============================================================

\begin{proposition}[Distortion Magnitude]
\label{prop:distortion}
The equilibrium adoption share of protocol $A$ among early movers, denoted $\sigma_A^*$, differs from the full-information optimal share $\sigma_A^{\text{FI}}$ by
\[
|\sigma_A^* - \sigma_A^{\text{FI}}| \leq O\!\left(\sqrt{\frac{\log N_1}{N_0}}\right).
\]
\end{proposition}

\begin{proof}[Proof sketch]
The information externality operates through the influence function $\phi_P$. Early movers internalize that their choice affects the late-mover MAB's posterior. The distortion magnitude is bounded using the MAB regret bound (Theorem~\ref{thm:regret}) and the smoothness of the logit response. As $N_0$ grows, the distortion vanishes at rate $O(\sqrt{\log N_1 / N_0})$.
\end{proof}

% ============================================================
\section{Conclusion}
% ============================================================

We have formalized audit protocol adoption as a sequential-move game with
incomplete information, where early movers face strategic uncertainty and
late movers deploy a UCB-based MAB strategy. The analysis reveals a novel
information-externality channel that distorts equilibrium adoption away
from the full-information benchmark. The regret bound and uniqueness
condition provide actionable guidance for protocol designers and ecosystem
architects.

% ============================================================
\bibliographystyle{plainnat}
\begin{thebibliography}{20}

\bibitem{goldwasser1989knowledge}
S.~Goldwasser, S.~Micali, and C.~Rackoff.
\newblock ``The knowledge complexity of interactive proof systems,''
\newblock \emph{SIAM Journal on Computing}, 18(1):186--208, 1989.

\bibitem{costan2016intel}
V.~Costan and S.~Devadas.
\newblock ``Intel SGX explained,''
\newblock \emph{IACR Cryptology ePrint Archive}, 2016.

\bibitem{nakamoto2008bitcoin}
S.~Nakamoto.
\newblock ``Bitcoin: A peer-to-peer electronic cash system,''
\newblock 2008.

\bibitem{mcfadden1974conditional}
D.~McFadden.
\newblock ``Conditional logit analysis of qualitative choice behavior,''
\newblock in \emph{Frontiers in Econometrics}, 1974.

\bibitem{auer2002finite}
P.~Auer, N.~Cesa-Bianchi, and P.~Fischer.
\newblock ``Finite-time analysis of the multiarmed bandit problem,''
\newblock \emph{Machine Learning}, 47(2):235--256, 2002.

\end{thebibliography}

\end{document}
